% Pillar-3 (group size) discussion + conclusion
\section{Discussion and Limitations}
\label{sec:p3-discussion}
Reframing $G$ as a preference-density dial explains why the zero-variance fraction
falls with $G$, but in our measured runs that fall did not translate into held-out
accuracy, so we resist naming an optimal $G$; our equivalence tests are powered to
reject large differences, not to certify exact parity. Read through the four-axis
pipeline factorization of \citet{ivison2024unpacking}, applied to verifiable-reward
stacks in the RLVR paradigm of \citet{lambert2024tulu3}, the same-stack audit
places the algorithm axis as residual ($\eta^2{=}0.023$; paired
$\Delta_{\text{GRPO}-\text{PPO}}{=}-0.002$, permutation $p{=}0.62$). On the 12-run
saturated $G$ sweep, $G$ accounts for $\eta^2{=}0.54$ of the variance in last-10
training reward and $0.63$ in gradient norm (Cohen's $d{=}-1.47$ for $G{=}2{\to}16$;
\texttt{experiments/results/berkeley/unpacking\_dpo\_ppo\_factorization.json}).
These are within-cohort variance shares, not accuracy effects: near the ceiling
the across-seed spread is so small that a $0.0066$ shift in mean reward is a large
standardized effect while staying inside the $\epsilon{=}0.02$ equivalence margin.
An independent Miller-style error-bars audit \citep{miller2024erroreval} agrees that
the native-Wu $G{=}2\approx G{=}16$ equivalence is borderline at $n{=}3$ (TOST
$p{=}0.056$ paired; held-out $p{=}0.100$), and its SNR-slope check is not
interpretable because the logged ``SNR'' is a transform of the
non-degenerate-step fraction (Section~\ref{sec:group-size-iter27})
(\texttt{experiments/results/berkeley/adding\_error\_bars\_summary.json}).

The scale of the evidence is the main limitation. The saturated sweep is $n{=}3$
seeds (minimum detectable effect $d{=}3.07$); the unsaturated iso-compute sweep is
a 0.5B model trained for $60$ steps with five seeds; the Qwen3-8B adapters have one
training seed per $G$ and fix steps rather than rollouts. None of these runs logs
gradient vectors, so whitening and leave-one-out effects on gradient variance are
inferred only through accuracy and ZVF. A multi-seed, matched-token sweep on
Qwen3-8B with MATH-style low-pass-rate problems, and the $2{\times}2$
contrast-versus-variance design, are the natural next experiments.

The all-pairs contrast identity is an exact algebraic rewriting whose observable
signatures are consistent with logged per-step diagnostics (ZVF and gradient
norms); direct gradient-vector validation remains future work. The
GRPO$\leftrightarrow$DPO bridge is exact only up to a KL penalty and the
projection of zero-variance groups; treating GRPO as literally DPO ignores
stabilizing structure the group baseline may carry.

\section{Conclusion}
\label{sec:p3-conclusion}
Group size in GRPO is best understood not as a variance knob to be maximized but as
a dial on how often a prompt yields any implicit preference pair. Measured across a
saturated 0.5B sweep, an unsaturated iso-compute 0.5B sweep up to $G{=}32$, and
four Qwen3-8B/GSM8K adapters, larger groups consistently lowered the zero-variance
fraction---whose measured level exceeds the homogeneous closed form because prompt
difficulty varies---but did not raise held-out accuracy, and removing GRPO's
$1/\sigma_r$ whitening changed nothing measurable. On this evidence the cheap end of
the range ($G\approx2$--$8$) loses nothing at equal compute in the settings tested;
whether larger groups pay off on harder, longer-horizon training remains open. We
release the sweeps, per-step logs, Colab driver and analysis scripts
(\texttt{platform\_colab/p3\_isocompute\_*.py},
\texttt{platform\_modal/scripts/group\_size\_review\_1ed8\_audit.py}).
